% ===== Джерела до розділу 9. Ключі з префіксом ch09: =====
% Позиції ядра (bib-core.tex) тут НЕ дублюються.

% --- Огляди та основи ML у безпеці ---
\bibitem{ch09:ucci} Ucci D., Aniello L., Baldoni R. Survey of Machine Learning Techniques for Malware Analysis. \emph{Computers \& Security}. 2019. Vol.~81. P.~123--147. DOI: 10.1016/j.cose.2018.11.001.

\bibitem{ch09:gibert} Gibert D., Mateu C., Planes J. The Rise of Machine Learning for Detection and Classification of Malware: Research Developments, Trends and Challenges. \emph{Journal of Network and Computer Applications}. 2020. Vol.~153. Art.~102526.

\bibitem{ch09:sommer} Sommer R., Paxson V. Outside the Closed World: On Using Machine Learning for Network Intrusion Detection. \emph{2010 IEEE Symposium on Security and Privacy}. Oakland: IEEE, 2010. P.~305--316.

% --- Ознаки та моделі ---
\bibitem{ch09:schultz} Schultz M. G., Eskin E., Zadok E., Stolfo S. J. Data Mining Methods for Detection of New Malicious Executables. \emph{2001 IEEE Symposium on Security and Privacy}. Oakland: IEEE, 2001. P.~38--49.

\bibitem{ch09:saxe} Saxe J., Berlin K. Deep Neural Network Based Malware Detection Using Two Dimensional Binary Program Features. \emph{2015 10th International Conference on Malicious and Unwanted Software (MALWARE)}. Fajardo: IEEE, 2015. P.~11--20.

\bibitem{ch09:raff} Raff E., Barker J., Sylvester J., Brandon R., Catanzaro B., Nicholas C. Malware Detection by Eating a Whole EXE. \emph{AAAI Workshop on Artificial Intelligence for Cyber Security}. 2018.

\bibitem{ch09:ember} Anderson H. S., Roth P. EMBER: An Open Dataset for Training Static PE Malware Machine Learning Models. \emph{arXiv preprint arXiv:1804.04637}. 2018.

% --- Змагальне машинне навчання ---
\bibitem{ch09:biggio} Biggio B., Roli F. Wild Patterns: Ten Years After the Rise of Adversarial Machine Learning. \emph{Pattern Recognition}. 2018. Vol.~84. P.~317--331.

\bibitem{ch09:grosse} Grosse K., Papernot N., Manoharan P., Backes M., McDaniel P. Adversarial Examples for Malware Detection. \emph{Computer Security --- ESORICS 2017}. LNCS Vol.~10493. Cham: Springer, 2017. P.~62--79.

\bibitem{ch09:anderson-gym} Anderson H. S., Kharkar A., Filar B., Evans D., Roth P. Learning to Evade Static PE Machine Learning Malware Models via Reinforcement Learning. \emph{arXiv preprint arXiv:1801.08917}. 2018.

% --- Українські джерела ---
\bibitem{ch09:ml-ua} Застосування методів машинного навчання для виявлення шкідливого програмного забезпечення. \emph{Кібербезпека: освіта, наука, техніка}. 2024. URL: \url{https://csecurity.kubg.edu.ua/index.php/journal/article/view/688}.

\bibitem{ch09:neural-ua} Нейромережеві моделі класифікації шкідливого програмного забезпечення за статичними ознаками. \emph{Захист інформації}. 2024. Т.~26, №~1.
